\chapter{Plataforma Web para Exploração de Algoritmos de Regressão Simbólica}

\section{Panorama Geral}
A plataforma proposta nesse projeto tem como objetivo principal a manipulação de modelos de regressão simbólica em \textit{pipelines} de extração e processamento de dados. A plataforma pode ser separada em três principais componentes:
\begin{enumerate}
    \item Backend: Servidor responsável por processar as requisições Web, gerencias o banco de dados e os modelos de regressão simbólica;
    \item Worker: Programa especializado para treinamento dos modelos de regressão; e
    \item Frontend: Aplicação web que permite o usuário interagir com os modelos criados e as demais entidades do projeto.
\end{enumerate}

Para entender melhor as funcionalidades esperadas da plataforma, primeiro foram elencados os requisitos funcionais e não funcionais do projeto. Segundo o DDD, os requisitos funcionais são aqueles que dizem respeito a ações que podem ser realizadas dentro do sistema desenvolvido. Em outras palavras, os requisitos funcionais dizem o que uma aplicação é capaz de fazer. Por outro lado, os requisitos não funcionais dizem como um sistema opera, delimitando aspectos de segurança, restrições e desempenho da aplicação.

\section{Histórias de Usuário}
Como os requisitos funcionais dizem respeito às ações que podem ser realizadas dentro de uma aplicação, eles foram descritos por meio de Histórias de Usuário (\textit{User Stories}), técnica amplamente utilizada em metodologias ágeis \cite{cohn2004user}. Essas histórias descrevem um fluxo de operação do sistema sob a perspectiva do ator final, permitindo entender as funcionalidades e o valor de cada entrega.

Além disso, como um dos objetivos do projeto é a criação de uma plataforma de manipulação de algoritmos de regressão simbólica que tenha paridade com as ferramentas comerciais disponíveis (como TuringBot e HeuristicLab), incorporou-se funcionalidades consolidadas no mercado juntamente com as inovações propostas pelo Spectrum (como a linguagem IQL). Dessa forma, o projeto aqui descrito possui os seguintes requisitos funcionais, estratificados pelas seguintes histórias de usuário principais:

\begin{itemize}
    \item \textbf{US1 - Upload de Dados}: Como pesquisador, quero fazer \textit{upload} de arquivos CSV para que eu possa utilizar meus próprios dados experimentais no treinamento de modelos.
    \item \textbf{US2 - Datasets Aleatórios}: Como desenvolvedor, quero gerar conjuntos de dados aleatórios para testar o comportamento dos algoritmos em diferentes cenários matemáticos.
    \item \textbf{US3 - Treinamento de Modelos}: Como usuário, quero configurar os meta-parâmetros do algoritmo Eggp para que eu possa otimizar a busca por expressões que respeitem restrições específicas de domínio.
    \item \textbf{US4 - Consulta IQL}: Como analista de dados, quero consultar expressões encontradas através da linguagem IQL para identificar padrões matemáticos de interesse.
    \item \textbf{US5 - Visualização de Métricas}: Como pesquisador, quero observar métricas de performance (MSE, R²) e complexidade estrutural para facilitar a seleção do modelo mais adequado.
    \item \textbf{US6 - Predição e Validação}: Como usuário, quero utilizar os modelos treinados para realizar previsões em novos dados, validando a generalização da solução encontrada.
\end{itemize}

\section{Interface Web e Blocos de Visualização}
A plataforma Spectrum utiliza uma interface baseada em "blocos" para visualização e manipulação de dados, permitindo que a interface se adapte às necessidades de cada perfil de usuário. Cada perfil pode ser personalizado com blocos modulares, como por exemplo:
\begin{itemize}
    \item \textbf{Bloco de Editor IQL}: Provê uma interface de codificação com realce de sintaxe para execução de consultas declarativas sobre os modelos de regressão;
    \item \textbf{Bloco de Markdown}: Permite a documentação textual diretamente na plataforma, integrando anotações do pesquisador aos experimentos realizados.
\end{itemize}

A interface em blocs se assemelha a outras plataformas amplamente utilizadas em ambientes de ciência de dados, como o Google Colab. Dessa forma, a plataforma proposta se aproxima de outras ferramentas já utilizadas atualmente, facilitando a utilização dela. Além disso, o objetivo dessa interface modular é integrar as etapas de exploração, escrita e utilização dos modelos de regressão. Ferramentas disponíveis hoje em dia, como HeuristicLab e TuringBot, não permitem a documentação do modelo de forma integrada~\cite{heuristiclabDocumentationAboutHeuristicLabx2013,turingbot2020}, implicando na necessidade de um pesquisador utilizar diversas ferramentas simultaneamente e acarretando em trocas de contexto constantes.

\section{Requisitos não-funcionais}
Os requisitos não-funcionais definem as restrições e qualidades do sistema Spectrum. Para esta plataforma, os seguintes requisitos foram estabelecidos:
\begin{enumerate}
    \item \textbf{Segurança}: A autenticação deve ser realizada via JSON Web Tokens (JWT) armazenados em cookies \textit{HttpOnly} para mitigar ataques de Cross-Site Scripting (XSS);
    \item \textbf{Desempenho}: O treinamento de modelos, por ser uma tarefa intensiva em termos de CPU, deve ser executado de forma assíncrona em processos separados (\textit{workers});
    \item \textbf{Integridade}: Todas as mutações no banco de dados devem seguir o padrão \textit{Unit of Work}, garantindo a atomicidade das transações;
    \item \textbf{Portabilidade}: O sistema deve ser capaz de ser executado em diferentes ambientes por meio de conteinerização com Docker.
\end{enumerate}

\section{Back-End}
O back-end do projeto foi desenvolvido seguindo o padrão monólito modularizado, em oposição a uma arquitetura de micro-serviços pura, visando reduzir a complexidade de rede sem abrir mão da separação de responsabilidades. O servidor é responsável por gerenciar as principais entidades do sistema:
\begin{itemize}
    \item \textbf{User}: Representa os usuários da plataforma;
    \item \textbf{Dataset}: Metadados e referências aos arquivos CSV de dados;
    \item \textbf{Job}: Representa uma configuração para o algoritmo eggp, contendo os meta parâmetros do algoritmo;
    \item \textbf{Run}: O resultado de um treinamento, contendo metadados como tempo de execução;
    \item \textbf{Profile}: Personalização da interface e blocos de visualização para o usuário;
    \item \textbf{Model}: Representa um modelo gerado por um \textbf{Job}, com referências para o arquivo egraph gerado;
    \item \textbf{Block}: Representa um bloco no editor;
    \item \textbf{Inference Run}: Representa uma requisição em linguagem IQL para determinada \textbf{Profile};
    \item \textbf{Inference Result}: Representa os resultados de uma consulta IQL.
\end{itemize}

As sub-seções seguintes apresentam as rotas disponibilizadas por cada módulo para manipulação do seu domínio, bem como os valores esperadas de entrada e saída para cada rota.

\subsection{Usuários}
Como forma de permitir a criação de modelos privados, a plataforma Spectrum oferece uma funcionalidade de criação de contas de usuário. Todas as demais entidades do sistem estão relacionadas, direta ou indiretamente, com um usuário. Usuários representam pessoas que utilizam a plataforma. Um usuário é definido a nível de banco de dados pelo código \ref{code:user_db}, e a nível de domínio pelo código \ref{code:user_domain}.

\begin{lstlisting}[language=Sql, caption={Código SQL para criação da tabela de usuários}, label={code:user_db}]
SQL CODE
\end{lstlisting}

\begin{lstlisting}[language=Python, caption={Classe de domínio que define um usuário}, label={code:user_domain}]
from datetime import datetime
from typing import Optional

from pydantic import Field, EmailStr, SecretStr
from ..base import BaseMutableDomainModel


class User(BaseMutableDomainModel):
    first_name: str = Field(..., description="The user's first name")

    last_name: str = Field(..., description="The user's last name")

    email: EmailStr = Field(..., description="The user's email address")

    password_hash: SecretStr = Field(...,
                                     description="The hashed password of the user")

    deleted_at: Optional[datetime] = Field(
        None, description="The timestamp when the user was deleted, if applicable")

    @classmethod
    def new(
        cls,
        *,
        first_name: str,
        last_name: str,
        email: EmailStr,
        password_hash: SecretStr,
    ):
        return cls(
            first_name=first_name,
            last_name=last_name,
            email=email,
            password_hash=password_hash,
            deleted_at=None,
        )
\end{lstlisting}

\section{Workers}

\section{Front-End}

